\section{LLM Fallback Details}
\label{app:fallback}

In the fallback, the LLM acts as a semantic summarization component for the missing symbolic‑execution step. It receives the current code fragment, available symbolic states and path conditions, previously generated callee specifications or loop summaries, and an ACSL skeleton. The LLM is tasked with following the relevant paths, tracking changes to variables and memory locations, and proposing ACSL clauses that summarize the observed behavior. The resulting candidates must still be proved by Frama‑C/WP. 

%The fallback uses the LLM as a semantic summarization component for the missing symbolic-execution step, proposing ACSL candidates that must still be proved by Frama-C/WP. It receives the current code fragment, available symbolic states and path conditions, previously generated callee specifications or loop summaries, and an ACSL skeleton, and is instructed to follow the relevant paths, to track changes to variables and memory locations, and to propose ACSL clauses that summarize the observed behavior.

We use this fallback for four function-level tasks. \emph{Precondition generation} asks the LLM to infer candidate \texttt{requires} clauses describing the validity, separation and bound constraints. \emph{Assigns-clause generation} asks the LLM to infer the function's footprint, i.e., \texttt{assigns} clauses describing which visible memory locations may be modified. \emph{Postcondition generation} asks the LLM to summarize the relation between the pre-state and the post-state, filling candidate \texttt{ensures} clauses. \emph{Inductive-predicate generation} asks the LLM to decide whether the program requires an inductive predicate and, when needed, to propose one for capturing inductive relations in recursive data structures or recursive functions.

